\chapter{Build and Command Line}

FEMaster is built with CMake. A typical build produces an executable named \file{FEMaster} in the build output directory. The executable reads an input deck and writes result files using a common job base name. If no explicit output path is supplied, the job base is derived from the input file name. FEMaster then writes both the native text result file \file{.res} and the CalculiX/CGX result file \file{.frd}.

\section{Build}

\begin{shellcmd}
cmake -S . -B build -DCMAKE_BUILD_TYPE=Release
cmake --build build --config Release
\end{shellcmd}

A release build is the normal choice for analysis work. A debug build is useful for small diagnostic models, assertions, and development checks, but it is usually too slow for larger finite-element models.

\section{CMake presets}

The repository provides CMake presets for the common supported toolchain combinations. Presets place their build directory under \file{build/<preset-name>} and set the normal release-oriented defaults: tests disabled, double precision enabled, Eigen compile-time reduction flags enabled, and OpenMP disabled unless the preset explicitly enables it.

\begin{shellcmd}
cmake --preset linux
cmake --build --preset linux

cmake --preset linux-mkl-omp
cmake --build --preset linux-mkl-omp
\end{shellcmd}

The Linux presets cover CPU-only builds, OpenMP builds, oneMKL builds, CUDA builds, cuDSS builds, and combined CUDA/oneMKL configurations. Windows presets are split between MSVC and IntelLLVM toolchains. macOS currently provides CPU presets. The exact preset list should be read from \file{CMakePresets.json} for the checked-out version.

\begin{longtable}{@{}p{0.32\linewidth}p{0.58\linewidth}@{}}
\toprule
\textbf{CMake option} & \textbf{Meaning} \\
\midrule
\endhead
\key{FEMASTER_BUILD_TESTS} & Build and register the GoogleTest executable. \\
\key{FEMASTER_ENABLE_CUDA} & Build CUDA solver backends. \\
\key{FEMASTER_ENABLE_CUDSS} & Use NVIDIA cuDSS for CUDA sparse direct solves. Requires CUDA. \\
\key{FEMASTER_ENABLE_MKL} & Use Intel oneMKL for Eigen operations and PARDISO. \\
\key{FEMASTER_MKL_SEQUENTIAL} & Use the sequential oneMKL threading layer. \\
\key{FEMASTER_MKL_STATIC} & Link oneMKL component libraries statically. \\
\key{FEMASTER_ENABLE_OPENMP} & Enable OpenMP for FEMaster CPU loops. \\
\key{FEMASTER_DOUBLE_PRECISION} & Use double precision arithmetic. \\
\key{FEMASTER_FAST_EIGEN_COMPILE} & Use Eigen compile-time reduction flags. \\
\key{FEMASTER_STATIC_RUNTIME} & Link compiler C/C++ runtimes statically where supported. \\
\key{FEMASTER_FULLY_STATIC} & Produce a fully static Linux executable. Incompatible with CUDA shared libraries. \\
\key{FEMASTER_TIME_REPORT} & Enable compiler timing diagnostics where supported. \\
\key{FEMASTER_DEPS_INCLUDE_DIR} & Include root containing header-only dependencies such as Eigen, Spectra, and argparse. \\
\key{FEMASTER_CUDSS_ROOT} & Root of an NVIDIA cuDSS installation. Defaults to the \key{CUDSS_DIR} environment variable. \\
\bottomrule
\end{longtable}

\section{Running an input deck}

\begin{shellcmd}
FEMaster model.inp
FEMaster model.inp --output model.res
FEMaster model.inp --ncpus 8
\end{shellcmd}

\begin{longtable}{@{}p{0.28\linewidth}p{0.62\linewidth}@{}}
\toprule
\textbf{Argument} & \textbf{Meaning} \\
\midrule
\endhead
\file{input_file} & Path to the input deck. If no \file{.inp} suffix is supplied, it is added automatically. \\
\key{--output FILE} & Explicit output base or result file path. The suffix \file{.res}, \file{.frd}, or \file{.inp} is stripped before writers are opened, so \file{--output run.res} produces \file{run.res} and \file{run.frd}. \\
\key{--ncpus N} & Number of threads allowed for parallel assembly and compute sections. \\
\bottomrule
\end{longtable}

At the end of a regular solver run, FEMaster prints a process timing summary in hours, minutes, seconds, milliseconds, and total milliseconds. This timing covers the full executable run, including parsing, solving, result writing, and shutdown.

\section{DSL documentation mode}

The executable can print DSL registration information for the installed version. This is useful when checking which commands, keywords, and variants a particular build accepts.

\begin{shellcmd}
FEMaster --document --doc-list
FEMaster --document --doc-show CLOAD
FEMaster --document --doc-variants ELASTIC
\end{shellcmd}

\begin{longtable}{@{}p{0.34\linewidth}p{0.56\linewidth}@{}}
\toprule
\textbf{Option} & \textbf{Meaning} \\
\midrule
\endhead
\key{--doc-list} & Print the available DSL commands. \\
\key{--doc-show CMD} & Show details for one command. \\
\key{--doc-tokens CMD} & Show expected tokens for a command. \\
\key{--doc-variants CMD} & Show accepted data-line variants. \\
\key{--doc-search TEXT} & Search command names and descriptions. \\
\key{--doc-where-token TOKEN} & Find commands that use a token. \\
\key{--doc-all} & Print the full command overview. \\
\key{--doc-format text|md|json} & Select the documentation output format. Currently the implementation falls back to text for non-text formats. \\
\key{--doc-verbosity index|compact|full} & Select the detail level for command documentation. \\
\key{--doc-width N} & Set the wrapping width for documentation output. Use 0 for no wrapping. \\
\key{--doc-no-wrap} & Disable wrapping. Equivalent to \key{--doc-width 0}. \\
\key{--doc-regex} & Treat \key{--doc-search} as a regular expression. \\
\bottomrule
\end{longtable}
